\subsection{Primary Decomposition and Cyclic Primary Modules}

\begin{definition} \label{def:primary-module-over-pid} \hlt{Primary Module}\\
Let $p\in R$ be prime. A torsion $R$-module $M$ is \hlt{$p$-primary} if $(\exists e\in\N)\,[\operatorname{Ann}_R(M)=\langle p^e\rangle]$.\\
Equivalently, $(\exists e\in\N)\,[o(M)\sim p^e]$ by Definition \ref{def:module-order-over-pid}.
\end{definition}

\begin{theorem} \label{thm:primary-decomposition-pid-modules} \hlt{Primary Decomposition}\\
Let $n\in\N$, $M\neq\{0_M\}$ be f.g. torsion, $p_1,\ldots,p_n\in R$ be distinct nonassociate primes. Suppose $e_1,\ldots,e_n\in\N$ and $o(M)\sim\mu\defeq\prod\limits_{i=1}^np_i^{e_i}$.\\
Set $I\defeq\{1,\ldots,n\}$. For $i\in I$, set $q_i\defeq p_i^{e_i}$, $\mu_i\defeq\prod\limits_{j\in I\setminus\{i\}}q_j$, and $M_{p_i}\defeq\mu_iM$.
\begin{enumerate}[label=(\roman*)]
\setlength{\itemsep}{0pt}
\item $M=\bigoplus\limits_{i=1}^nM_{p_i}$ and $(\forall i\in I)\,[M_{p_i}=\{v\in M:q_iv=0_M\}]$.
\item Each $M_{p_i}$ is f.g. and $p_i$-primary, with $o(M_{p_i})\sim q_i$.
\item The unordered family $\{(p_i,M_{p_i}):i\in I\}$, with primes taken up to associates, is determined by $M$.
\end{enumerate}
\begin{proof}
\begin{enumerate}[label=(\roman*)]
\setlength{\itemsep}{0pt}
\item
\begin{enumerate}[label=(\alph*)]
\setlength{\itemsep}{0pt}
\item \hlt{Spanning}: The elements $\mu_1,\ldots,\mu_n$ have no common nonunit divisor by the distinct-prime factorization of $\mu$.\\
Choose $a_1,\ldots,a_n\in R$ with $\sum\limits_{i=1}^na_i\mu_i=1_R$ by induction on $n$ from Theorem \ref{thm:pid-prime-irreducible-bezout}(iii).\\
Thus $(\forall x\in M)\,[x=\sum\limits_{i=1}^na_i\mu_ix]$, so $M=\sum\limits_{i=1}^nM_{p_i}$.
\item \hlt{Kernel Description}: Since $q_i\mu_i=\mu$ and $\mu M=\{0_M\}$, $(\forall i\in I)\,[q_iM_{p_i}=\{0_M\}]$.\\
Fix $i\in I$ and $v\in M$ with $q_iv=0_M$. The elements $\mu_i$ and $q_i$ have no common nonunit divisor.\\
Choose $a,b\in R$ with $a\mu_i+bq_i=1_R$ by Theorem \ref{thm:pid-prime-irreducible-bezout}(iii).\\
Then $v=a\mu_iv+bq_iv=a\mu_iv\in M_{p_i}$. Hence $M_{p_i}=\{v\in M:q_iv=0_M\}$.
\item \hlt{Directness}: Let $(x_i)_{i=1}^n\in\prod\limits_{i=1}^nM_{p_i}$ satisfy $\sum\limits_{i=1}^nx_i=0_M$, and fix $i\in I$.\\
If $j\neq i$, then $q_j\mid\mu_i$ and $q_jx_j=0_M$, so $\mu_ix_j=0_M$. By multiplying the displayed sum by $\mu_i$, $\mu_ix_i=0_M$.\\
Together with $q_ix_i=0_M$ and the Bézout identity for $\mu_i,q_i$, this implies $x_i=0_M$. Thus the sum is direct.
\end{enumerate}
\item The image $M_{p_i}=\mu_iM$ is f.g. because the $\mu_i$-multiples of a finite generating set of $M$ generate it. For $i\in I$, let $\alpha_i\defeq o(M_{p_i})$.\\
Since $q_iM_{p_i}=\{0_M\}$, $(\forall i\in I)\,[\alpha_i\mid q_i]$. By Theorem \ref{thm:pid-is-ufd}(i)--(ii), choose $f_i\in\{0,\ldots,e_i\}$ with $\alpha_i\sim p_i^{f_i}$ $(i\in I)$.\\
By part (i), Proposition \ref{prop:orders-direct-sum-lcm}, $o(M)\sim\operatorname{lcm}(\alpha_1,\ldots,\alpha_n)$. Hence $\operatorname{lcm}(\alpha_1,\ldots,\alpha_n)\sim\prod\limits_{i=1}^nq_i$.\\
Unique factorization forces $(\forall i\in I)\,[f_i=e_i]$. Since $q_i\neq0_R$ and $q_iM_{p_i}=\{0_M\}$, the module $M_{p_i}$ is torsion.\\
Thus $o(M_{p_i})\sim q_i$, so $M_{p_i}$ is $p_i$-primary by Definition \ref{def:primary-module-over-pid}.
\item The prime-power factors $q_i$ of $o(M)$ are determined up to associates and permutation by Theorem \ref{thm:pid-is-ufd}(ii).\\
If $q_i'\sim q_i$, then $\{v\in M:q_i'v=0_M\}=\{v\in M:q_iv=0_M\}$ because the associate factor is a unit.\\
Part (i) therefore makes each primary component intrinsic.
\end{enumerate}
\noindent
\hlt{Structural Commentary}: Bézout isolates prime-power kernels; their order lcm recovers every primary exponent.\\
\hlt{Key Insight}: Each intrinsic component is exactly $\ker(q_i:M\to M)$.
\end{proof}
\end{theorem}

\begin{definition} \label{def:p-power-and-p-torsion-submodules} \hlt{$p$-Power and $p$-Torsion Submodules}\\
Let $p\in R$ be prime and let $M$ be an $R$-module.
\begin{enumerate}[label=(\roman*)]
\setlength{\itemsep}{0pt}
\item For $j\in\N_0$, the \hlt{$j$th $p$-power submodule} is $p^jM\defeq\{p^jv:v\in M\}=\operatorname{im}(p^j\operatorname{id}_M)\leq M$.
\item The \hlt{first $p$-torsion layer} is $M[p]\defeq\ker(p\operatorname{id}_M)=\{v\in M:pv=0_M\}\leq M$.
\end{enumerate}
\end{definition}

\begin{lemma} \label{lem:p-torsion-layer-vector-space} \hlt{The First $p$-Torsion Layer}\\
Let $p\in R$ be prime, let $\kappa\defeq R/\langle p\rangle$, and let $M$ be an $R$-module.
\begin{enumerate}[label=(\roman*)]
\setlength{\itemsep}{0pt}
\item If $N$ is an $R$-module and $pN=\{0_N\}$, then $(r+\langle p\rangle)n\defeq rn$ $(r\in R,n\in N)$ defines a $\kappa$-vector-space structure on $N$.\\
In particular, $M[p]$ is a $\kappa$-vector space.
\item If $M=S\oplus T$, then $M[p]=S[p]\oplus T[p]$.
\end{enumerate}
\begin{proof}
\begin{enumerate}[label=(\roman*)]
\setlength{\itemsep}{0pt}
\item The prime element $p$ is irreducible, so $\langle p\rangle$ is maximal by Theorem \ref{thm:pid-prime-irreducible-bezout}(i). Hence $\kappa$ is a field by Theorem \ref{thm:quotient-by-maximal-ideal-is-field}.\\
Suppose $r+\langle p\rangle=r'+\langle p\rangle$ and fix $n\in N$. Choose $c\in R$ satisfying $r-r'=cp$. Then $(r-r')n=c(pn)=0_N$.\\
Thus the scalar action is well-defined. The vector-space axioms follow from the $R$-module axioms and the quotient-ring operations.\\
The submodule $M[p]$ satisfies $p(M[p])=\{0_M\}$ by Definition \ref{def:p-power-and-p-torsion-submodules}(ii), so the particular case follows.
\item Let $M=S\oplus T$ and fix $v\in M[p]$. Write $v=s+t$ with $(s,t)\in S\times T$. Then $ps+pt=0_M$.\\
Since $ps\in S$, $pt\in T$, and $S\cap T=\{0_M\}$, one has $ps=pt=0_M$. Hence $v\in S[p]+T[p]$.\\
Conversely, $s\in S[p]$ and $t\in T[p]$ imply $p(s+t)=0_M$. Finally, $S[p]\cap T[p]\subseteq S\cap T=\{0_M\}$, so the sum is direct.
\end{enumerate}
\hlt{Structural Commentary}: Annihilation by $p$ permits residue-class scalars; direct-sum intersections separate the $p$-killed coordinates.\\
\hlt{Key Insight}: The functor $M\mapsto M[p]$ preserves finite internal direct sums.
\end{proof}
\end{lemma}

\begin{definition} \label{def:p-primary-layer-dimensions} \hlt{Primary Layer Dimensions}\\
Let $p\in R$ be prime, let $\kappa\defeq R/\langle p\rangle$, and let $M$ be an $R$-module. For $k\in\N$, one has $p^kM\leq p^{k-1}M$ because $(\forall v\in M)\,[p^kv=p^{k-1}(pv)]$.\\
Every class $p^{k-1}v+p^kM$ is annihilated by $p$. Thus $p^{k-1}M/p^kM$ is a $\kappa$-vector space by Lemma \ref{lem:p-torsion-layer-vector-space}(i).\\
Define the \hlt{$k$th primary layer dimension} by $\lambda_k(M)\defeq\dim_{\kappa}(p^{k-1}M/p^kM)$.
\end{definition}

\begin{proposition} \label{prop:p-primary-layer-counts-recover-exponents} \hlt{Layer Counts Recover Primary Exponents}\\
Let $p\in R$ be prime, $\kappa\defeq R/\langle p\rangle$, $m\in\N$, $e_1\geq\cdots\geq e_m\geq1$. Suppose $M\cong\bigoplus\limits_{i=1}^mR/\langle p^{e_i}\rangle$.
\begin{enumerate}[label=(\roman*)]
\setlength{\itemsep}{0pt}
\item $(\forall k\in\N)\,[\lambda_k(M)=\#\{i\in\{1,\ldots,m\}:e_i\geq k\}]$.
\item $(\forall k\in\N)\,[\#\{i\in\{1,\ldots,m\}:e_i=k\}=\lambda_k(M)-\lambda_{k+1}(M)]$.
\end{enumerate}
\begin{proof}
\begin{enumerate}[label=(\roman*)]
\setlength{\itemsep}{0pt}
\item Fix $(e,k)\in\N^2$ and set $C_e\defeq R/\langle p^e\rangle$ and $L_{e,k}\defeq p^{k-1}C_e/p^kC_e$. If $k>e$, then $L_{e,k}=\{0_{L_{e,k}}\}$.\\
Assume $k\leq e$ and define $\phi_{e,k}:\kappa\to L_{e,k}$ by $\phi_{e,k}(r+\langle p\rangle)\defeq p^{k-1}(r+\langle p^e\rangle)+p^kC_e$.\\
If $r-r'=ps$, then the two images differ by $p^k(s+\langle p^e\rangle)\in p^kC_e$, so $\phi_{e,k}$ is well-defined.\\
The map $\phi_{e,k}$ is $\kappa$-linear and surjective by its formula. If $\phi_{e,k}(r+\langle p\rangle)=0$, choose $s,t\in R$ satisfying $p^{k-1}(r-ps)=p^et$.\\
By cancellation in the domain, $r=ps+p^{e-k+1}t\in\langle p\rangle$. Thus $\phi_{e,k}$ is injective, and consequently $L_{e,k}\cong\kappa$ for $k\leq e$.\\
An $R$-module isomorphism maps $p^j$-power submodules isomorphically, induces isomorphisms on corresponding primary layers.\\
Thus $\lambda_k$ is invariant under $R$-module isomorphism.\\
For the displayed cyclic direct sum, $p^j$-power submodules and their quotients are coordinatewise.\\
Set $I_k\defeq\{i\in\{1,\ldots,m\}:e_i\geq k\}$. Hence its $k$th primary layer is $\bigoplus\limits_{i\in I_k}\kappa$, so $\lambda_k(M)=\#I_k$.
\item Fix $k\in\N$ and set $I_k\defeq\{i\in\{1,\ldots,m\}:e_i\geq k\}$. Then $I_{k+1}\subseteq I_k$ and $I_k\setminus I_{k+1}=\{i\in\{1,\ldots,m\}:e_i=k\}$.\\
By part (i), $\#(I_k\setminus I_{k+1})=\#I_k-\#I_{k+1}=\lambda_k(M)-\lambda_{k+1}(M)$.
\end{enumerate}
\noindent
\hlt{Structural Commentary}: A cyclic block contributes one residue-field dimension until its exponent is exhausted.\\
\hlt{Key Insight}: Consecutive layer-dimension differences count blocks ending at each exponent.
\end{proof}
\end{proposition}

\begin{theorem} \label{thm:cyclic-decomposition-primary-pid} \hlt{Cyclic Decomposition of a Primary Module}\\
Let $p\in R$ be prime, let $e\in\N$, and let $M$ be a f.g. $p$-primary module with $o(M)\sim p^e$.
\begin{enumerate}[label=(\roman*)]
\setlength{\itemsep}{0pt}
\item There are $m\in\N$, $v_1,\ldots,v_m\in M$, and $e=e_1\geq e_2\geq\cdots\geq e_m\geq1$. The vectors satisfy $M=\bigoplus\limits_{i=1}^m\langle\!\langle v_i\rangle\!\rangle_R$.\\
Their annihilators satisfy $(\forall i\in\{1,\ldots,m\})\,[\operatorname{Ann}_R(v_i)=\langle p^{e_i}\rangle]$.
\item The exponent multiset $\{e_1,\ldots,e_m\}$ is determined by $M$. Consequently, so is the ordered chain $\langle p^{e_1}\rangle\subseteq\cdots\subseteq\langle p^{e_m}\rangle$.
\end{enumerate}
\begin{proof}
\begin{enumerate}[label=(\roman*)]
\setlength{\itemsep}{0pt}
\item
\begin{enumerate}[label=(\alph*)]
\setlength{\itemsep}{0pt}
\item \hlt{Diagonal Presentation}: Choose $R^a\xrightarrow{A}R^b\to M\to0$ by Proposition \ref{prop:finitely-generated-pid-module-finite-presentation}. Choose Smith form $D$ of $A$ by Theorem \ref{thm:smith-normal-form-over-pid}.\\
Choose $s,t\in\N_0$ and nonunits $d_1,\ldots,d_t\in R$ with $M\cong R^s\oplus\bigoplus\limits_{\ell=1}^tR/\langle d_\ell\rangle$ by Corollary \ref{cor:smith-normal-form-presents-pid-modules}.\\
If $s>0$, the free summand $R^s$ contains a nonzero torsion-free element, contradicting that $M$ is torsion.\\
Hence $s=0$. Since $M\neq\{0_M\}$, one has $t\geq1$.
\item \hlt{Prime Powers}: The order identity $\operatorname{lcm}(d_1,\ldots,d_t)\sim o(M)\sim p^e$ holds by Proposition \ref{prop:orders-direct-sum-lcm}.\\
Every prime factor of $d_\ell$ is associate to $p$. By Theorem \ref{thm:pid-is-ufd}(i)--(ii), choose $e_\ell\in\N$ with $d_\ell\sim p^{e_\ell}$ $(1\leq\ell\leq t)$.\\
Set $m\defeq t$ and reorder so $e_1\geq\cdots\geq e_m$. The least-common-multiple identity implies $e_1=e$.
\item \hlt{Internal Summands}: Set $C\defeq\bigoplus\limits_{i=1}^mR/\langle p^{e_i}\rangle$. Since $\langle d_i\rangle=\langle p^{e_i}\rangle$, choose an isomorphism $\Psi:C\to M$.\\
Let $\mathbf{u}_i$ be the standard generator of the $i$th coordinate summand and set $v_i\defeq\Psi(\mathbf{u}_i)$.\\
Since $\Psi$ is an isomorphism, $M=\bigoplus\limits_{i=1}^m\langle\!\langle v_i\rangle\!\rangle_R$. The standard generator of $R/\langle p^{e_i}\rangle$ has annihilator $\langle p^{e_i}\rangle$.\\
An isomorphism preserves annihilators, so $\operatorname{Ann}_R(v_i)=\langle p^{e_i}\rangle$.
\end{enumerate}
\item The dimensions $\lambda_k(M)$ are intrinsic because module isomorphisms induce isomorphisms on the layer quotients.\\
By Proposition \ref{prop:p-primary-layer-counts-recover-exponents}(ii), the multiplicity of the exponent $k$ is $\lambda_k(M)-\lambda_{k+1}(M)$.\\
Thus $M$ determines every exponent multiplicity and hence the exponent multiset.\\
Its decreasing order determines the annihilator chain.
\end{enumerate}
\noindent
\hlt{Structural Commentary}: Smith form supplies prime-power blocks; primary-layer differences recover their multiplicities.\\
\hlt{Key Insight}: The differences $\lambda_k(M)-\lambda_{k+1}(M)$ recover each exponent multiplicity.
\end{proof}
\end{theorem}
